% math0 manuscript. Build with: latexmk -pdf paper.tex
\documentclass[11pt]{amsart}
\usepackage[margin=1.1in]{geometry}
\usepackage{amsmath,amssymb,amsthm,mathtools}
\usepackage{url}
\usepackage{hyperref}
\usepackage{xcolor}
\usepackage{listings}
\usepackage{newunicodechar}

\lstdefinelanguage{lean}{
  keywords={theorem,lemma,def,abbrev,example,fun,by,have,show,from,let,in,if,then,else,match,with,
    namespace,end,open,import,section,variable,instance,structure,where,calc,obtain,intro,exact,
    refine,rcases,rintro,simp,omega,decide,norm_num,linarith,nlinarith,positivity,induction,cases,
    private,inductive,deriving},
  sensitive=true, comment=[l]{--}, morecomment=[s]{/-}{-/}, string=[b]"}
\lstset{language=lean, basicstyle=\ttfamily\footnotesize, columns=fullflexible, keepspaces=true, breaklines=true,
  keywordstyle=\color{blue!60!black}\bfseries, commentstyle=\color{gray}, frame=single,
  extendedchars=true, inputencoding=utf8,
  literate={ℕ}{{$\mathbb{N}$}}1 {ℤ}{{$\mathbb{Z}$}}1 {∀}{{$\forall$}}1
    {↔}{{$\leftrightarrow$}}1 {¬}{{$\neg$}}1 {→}{{$\to$}}1 {≤}{{$\le$}}1 {∈}{{$\in$}}1
    {∧}{{$\wedge$}}1 {∨}{{$\vee$}}1 {∑}{{$\sum$}}1 {·}{{$\cdot$}}1}
% Unicode commonly found in Lean sources.
\newunicodechar{ℕ}{\ensuremath{\mathbb{N}}}
\newunicodechar{ℤ}{\ensuremath{\mathbb{Z}}}
\newunicodechar{ℚ}{\ensuremath{\mathbb{Q}}}
\newunicodechar{ℝ}{\ensuremath{\mathbb{R}}}
\newunicodechar{ℂ}{\ensuremath{\mathbb{C}}}
\newunicodechar{∀}{\ensuremath{\forall}}
\newunicodechar{∃}{\ensuremath{\exists}}
\newunicodechar{→}{\ensuremath{\to}}
\newunicodechar{←}{\ensuremath{\leftarrow}}
\newunicodechar{↔}{\ensuremath{\leftrightarrow}}
\newunicodechar{↦}{\ensuremath{\mapsto}}
\newunicodechar{¬}{\ensuremath{\neg}}
\newunicodechar{∧}{\ensuremath{\wedge}}
\newunicodechar{∨}{\ensuremath{\vee}}
\newunicodechar{≤}{\ensuremath{\le}}
\newunicodechar{≥}{\ensuremath{\ge}}
\newunicodechar{≠}{\ensuremath{\ne}}
\newunicodechar{∈}{\ensuremath{\in}}
\newunicodechar{∉}{\ensuremath{\notin}}
\newunicodechar{⊆}{\ensuremath{\subseteq}}
\newunicodechar{∩}{\ensuremath{\cap}}
\newunicodechar{∪}{\ensuremath{\cup}}
\newunicodechar{∑}{\ensuremath{\sum}}
\newunicodechar{∏}{\ensuremath{\prod}}
\newunicodechar{⟨}{\ensuremath{\langle}}
\newunicodechar{⟩}{\ensuremath{\rangle}}
\newunicodechar{·}{\ensuremath{\cdot}}
\newunicodechar{∘}{\ensuremath{\circ}}
\newunicodechar{⁻}{\ensuremath{^{-}}}
\newunicodechar{¹}{\ensuremath{^{1}}}
\newunicodechar{²}{\ensuremath{^{2}}}
\newunicodechar{₀}{\ensuremath{_{0}}}
\newunicodechar{₁}{\ensuremath{_{1}}}
\newunicodechar{₂}{\ensuremath{_{2}}}
\newunicodechar{α}{\ensuremath{\alpha}}
\newunicodechar{β}{\ensuremath{\beta}}
\newunicodechar{ε}{\ensuremath{\varepsilon}}
\newunicodechar{δ}{\ensuremath{\delta}}
\newunicodechar{π}{\ensuremath{\pi}}
\newunicodechar{σ}{\ensuremath{\sigma}}
\newunicodechar{φ}{\ensuremath{\varphi}}
\newunicodechar{∞}{\ensuremath{\infty}}
\newunicodechar{⌊}{\ensuremath{\lfloor}}
\newunicodechar{⌋}{\ensuremath{\rfloor}}
\newunicodechar{⌈}{\ensuremath{\lceil}}
\newunicodechar{⌉}{\ensuremath{\rceil}}
\newunicodechar{∣}{\ensuremath{\mid}}
\newunicodechar{≡}{\ensuremath{\equiv}}
\newunicodechar{×}{\ensuremath{\times}}
\newunicodechar{⧸}{\ensuremath{/}}
\newunicodechar{‖}{\ensuremath{\|}}
\newunicodechar{√}{\ensuremath{\surd}}
\newunicodechar{∅}{\ensuremath{\emptyset}}
\newunicodechar{⦃}{\ensuremath{\{\!|}}
\newunicodechar{⦄}{\ensuremath{|\!\}}}
\newunicodechar{𝓝}{\ensuremath{\mathcal{N}}}
\newunicodechar{≫}{\ensuremath{\gg}}
\newunicodechar{≪}{\ensuremath{\ll}}

\theoremstyle{plain}
\newtheorem{theorem}{Theorem}
\newtheorem{lemma}[theorem]{Lemma}
\newtheorem{proposition}[theorem]{Proposition}
\newtheorem{corollary}[theorem]{Corollary}
\theoremstyle{definition}
\newtheorem{definition}[theorem]{Definition}
\newtheorem{remark}[theorem]{Remark}

\emergencystretch=3em
\allowdisplaybreaks
% Lean identifiers in running text (breakable at dots and underscores).
\expandafter\def\expandafter\UrlBreaks\expandafter{\UrlBreaks\do\_\do\.}
\DeclareUrlCommand\decl{\def\UrlFont{\normalfont\ttfamily}}
\newcommand{\N}{\mathbb{N}}
\newcommand{\Z}{\mathbb{Z}}
\newcommand{\Av}{\mathcal{A}}
\newcommand{\ib}[1]{[\,#1\,]}
\newcommand{\lb}{\char123}
\newcommand{\rb}{\char125}

\title[Schulte's alternating quadratic identity for Schr\"oder paths]{Schr\"oder paths counted by peaks: a formal proof of Schulte's alternating quadratic identity (OEIS A060693)}
\author{math0}
\address{Produced by Claude (an AI model developed by Anthropic) in the math0 project}
\date{October 21, 2026}

\begin{document}
\begin{abstract}
Let $T(n,k)$ be the number of Schr\"oder paths from $(0,0)$ to $(2n,0)$ with exactly $k$ peaks (OEIS A060693).
We prove that $\sum_{k=0}^{n}(-1)^k\,T(n,k)\,(n+1-k)^2=1+n+n^2$ for every $n\ge0$, as conjectured by Werner Schulte in the OEIS entry in 2017, by computing the alternating falling-factorial moments of order $0,1,2$ of the peak statistic, for paths started at any height.
The identity is not new: it follows quickly from Harris's formula $T_n(x)=C_n(x+1)$ (2007) relating the row polynomials to the Narayana polynomials, and a Lean proof was merged into the \texttt{trureturing} repository in September 2026; ours is an independent formal proof.
It is verified in Lean~4 against the frozen, fidelity-reviewed statement that math0 formalized from the OEIS entry, in which $T(n,k)$ is defined by enumerating paths and checked by the kernel against all $56$ listed terms.
\end{abstract}
\maketitle

\section{Introduction}

A \emph{Schr\"oder path} of semilength $n$ is a lattice path from $(0,0)$ to $(2n,0)$ with steps $U=(1,1)$, $D=(1,-1)$ and $H=(2,0)$ that never goes below the $x$-axis; a \emph{peak} is a step $U$ immediately followed by a step $D$. The number $T(n,k)$ of Schr\"oder paths of semilength $n$ with exactly $k$ peaks ($0\le k\le n$) is the triangle A060693 of the OEIS \cite{A060693}, submitted by F.~Chapoton in 2001; the path description is a comment of Emeric Deutsch (2003), which became the name of the entry in 2017.

On January 11, 2017, Werner Schulte added to the entry the line
\begin{equation}\label{eq:schulte}
  \sum_{k=0}^{n}(-1)^k\,T(n,k)\,(n+1-k)^2=1+n+n^2 ,
\end{equation}
labelled ``Conjecture''; two days earlier he had added, without that label, the linear companion $\sum_{k=0}^n(-1)^k(1+x(n-k))T(n,k)=x+(1-x)\ib{n=0}$. In revision \#147 of the entry (October~4, 2026) the line \eqref{eq:schulte} is still labelled ``Conjecture''. We prove it for all $n\ge0$ (Theorem~\ref{thm:main}), with a Lean proof of the full statement.

\subsection*{Prior work and status}
\begin{itemize}
\item \emph{A short proof from known formulas.} The entry records Mitch Harris's formula $T_n(x)=C_n(x+1)$ (January 2007), where $T_n(x)=\sum_kT(n,k)x^k$ and $C_n(x)=\sum_{k=1}^{n}N(n,k)x^k$, with $N(n,k)$ the Narayana numbers A001263 \cite{A001263}, is the $n$-th Narayana row polynomial (the formula is meant for $n\ge1$); equivalently, for all $n\ge0$, $T(n,k)=\sum_j N(n,j)\binom jk$ with $N$ the Narayana triangle A090181 \cite{A090181}, which has $N(0,0)=1$ (Ph.~Deléham, May 2007). Together with $N(n,1)=1$ and $N(n,2)=\binom n2$ for $n\ge1$, this gives \eqref{eq:schulte} in a few lines (Remark~\ref{rem:narayana}, not formalized). In substance, the identity is a corollary of formulas that were in the entry before the conjecture was posted.
\item \emph{A prior formal proof.} A Lean proof of \eqref{eq:schulte}, the theorem \texttt{schulte\_a060693} in the file \path{SchroederPeakAlternatingQuadraticMoment.lean}, was merged into the \texttt{dev} branch of the GitHub repository \texttt{the-omega-institute/trureturing} on September~13, 2026, in pull request \#7436 by ElonSG \cite{TRU}, which supersedes pull request \#7402 and was pre-registered in issue \#7384. According to the math0 literature referee, who read the source but could not compile it (nor have we), $T$ is defined there through a first-return decomposition of Schr\"oder words, and the proof uses recurrences for the signed falling-factorial moments of orders $0,1,2$. It was found by the math0 literature review after our proof was written; the two developments are independent.
\item The math0 referees found no other proof in arXiv, MathOverflow or Mathematics Stack Exchange.
\end{itemize}
We therefore classify our result as an \emph{independent formal proof} of a known result. Our proof uses only the combinatorial definition of $T$.

\section{Statement}\label{sec:statement}

The FORMULA section of A060693 \cite{A060693} (offset $0$, $0\le k\le n$) contains, verbatim, the line
\begin{quote}\raggedright
\texttt{Conjecture: Sum\_\lb k=0..n\rb\ (-1)\^{}k*T(n,k)*(n+1-k)\^{}2 = 1+n+n\^{}2. - \_Werner Schulte\_, Jan 11 2017}
\end{quote}

\begin{theorem}\label{thm:main}
For $n,k\ge0$ let $T(n,k)$ be the number of lattice paths from $(0,0)$ to $(2n,0)$ with steps $U=(1,1)$, $D=(1,-1)$, $H=(2,0)$ that never go below the $x$-axis and contain exactly $k$ occurrences of $U$ immediately followed by $D$. Then for every $n\ge0$,
\[
  \sum_{k=0}^{n}(-1)^k\,T(n,k)\,(n+1-k)^2=1+n+n^2 .
\]
\end{theorem}

\subsection*{Lean formalization}
The statement was formalized by math0 in \path{lean/Math0/Statements/OEIS/A060693.lean}, reviewed for fidelity, and frozen (SHA-256 \texttt{7877842d\ldots 6b90}) before any proof attempt. Its definitions and the statement read as follows (verbatim; we omit the import, comments, blank lines, all theorems, a \texttt{DecidablePred} instance for \decl{IsSchroderPath}, the definitions \decl{valid}, \decl{peaksFrom}, \decl{cnt}, \decl{dTerm}, \decl{step}, \decl{Tab}, \decl{look}, \decl{next}, \decl{tabs}, \decl{Tfast}, \decl{a} of the evaluator and data checks described below, and, at the start of the namespace, the inductive type \texttt{Step} with constructors \texttt{U}, \texttt{D}, \texttt{H}, its displacements \texttt{Step.dx : Step → ℕ} and \texttt{Step.dy : Step → ℤ}, giving $(1,1)$, $(1,-1)$, $(2,0)$, and the line \texttt{open Step}):
\begin{lstlisting}
namespace Math0.Statements.A060693
def width (p : List Step) : ℕ := (p.map Step.dx).sum
def height (p : List Step) : ℤ := (p.map Step.dy).sum
def IsSchroderPath (n : ℕ) (p : List Step) : Prop :=
  width p = 2 * n ∧ height p = 0 ∧ ∀ q ∈ p.inits, 0 ≤ height q
def peaks (p : List Step) : ℕ := (p.zip p.tail).count (U, D)
def words : ℕ → List (List Step)
  | 0 => [[]]
  | 1 => [[U], [D]]
  | m + 2 => (words (m + 1)).map (U :: ·) ++ (words (m + 1)).map (D :: ·) ++
      (words m).map (H :: ·)
def schroderPaths (n : ℕ) : Finset (List Step) :=
  (words (2 * n)).toFinset.filter (IsSchroderPath n)
def T (n k : ℕ) : ℕ := ((schroderPaths n).filter (fun p => peaks p = k)).card
def statement : Prop :=
  ∀ n : ℕ, ∑ k ∈ Finset.range (n + 1), (-1 : ℤ) ^ k * (T n k : ℤ) * ((n : ℤ) + 1 - k) ^ 2 =
    1 + (n : ℤ) + (n : ℤ) ^ 2
end Math0.Statements.A060693
\end{lstlisting}
A path is a word \texttt{p : List Step} drawn from the origin. For a word $p$ let $|p|=\#U+\#D+2\#H$ be its \emph{width} and $\operatorname{ht}(p)=\#U-\#D$ its \emph{height}; after a prefix $q$ the path is at $(|q|,\operatorname{ht}(q))$. \texttt{IsSchroderPath n p} says that the path ends at $(2n,0)$ and that every prefix (\texttt{p.inits} includes the empty prefix and $p$) ends at height $\ge0$; checking vertices suffices, since the path is straight between them. \texttt{peaks p} counts adjacent pairs $(U,D)$. The list \texttt{words m} contains every word of width $m$ exactly once (\decl{mem_words}, \decl{nodup_words}), so \texttt{T n k} is $T(n,k)$; the file proves \decl{T_eq_ncard}: \texttt{T n k = \lb p : List Step | IsSchroderPath n p ∧ peaks p = k\rb.ncard}. The sum is computed in $\Z$, so $n+1-k$ is not truncated, and \decl{statement} covers all $n\ge0$ with no further hypothesis.

\emph{Data checks.} As \texttt{T} enumerates exponentially many words, the file also defines a table-based evaluator \decl{Tfast} from the recursion of Lemma~\ref{lem:cnt}(b) below and proves \decl{T_eq_Tfast}: \texttt{T n k = Tfast n k} for all $n,k$. Through this lemma, \texttt{decide +kernel} checks in the kernel that \texttt{T} reproduces rows $n=0,\dots,9$ and $T(10,0)=16796$ (\decl{T_data}), and that the triangle read by rows reproduces all $56$ terms of the OEIS data (\decl{a_data}). This anchors the definition to the entry independently of the conjecture. The file also checks \eqref{eq:schulte} for $n\le12$ (\decl{statement_check}), as a transcription check only.

\section{Proof}\label{sec:proof}

The proof follows \path{lean/Math0/Frontier/OEIS/A060693.lean}. Throughout, $b$ ranges over $\{0,1\}$ (the Lean \texttt{false}, \texttt{true}), $\ib{P}$ is $1$ if $P$ holds and $0$ otherwise, and $\Av_m$ is the set of words of width $m$.

\subsection{Generalized counts}
The statement file defines, by recursion on words, \decl{valid} and \decl{peaksFrom}:
\[
\begin{aligned}
&\operatorname{valid}_h(\varepsilon)=\ib{h=0},\quad \operatorname{valid}_h(Up)=\operatorname{valid}_{h+1}(p),\quad \operatorname{valid}_0(Dp)=0,\\
&\operatorname{valid}_{h+1}(Dp)=\operatorname{valid}_h(p),\quad \operatorname{valid}_h(Hp)=\operatorname{valid}_h(p);\\
&\operatorname{pk}_b(\varepsilon)=0,\quad \operatorname{pk}_b(Up)=\operatorname{pk}_1(p),\quad \operatorname{pk}_b(Dp)=b+\operatorname{pk}_0(p),\quad \operatorname{pk}_b(Hp)=\operatorname{pk}_0(p).
\end{aligned}
\]
So $\operatorname{valid}_h(p)$ holds if and only if the path $p$ started at height $h$ never goes below the axis and ends on it (\decl{valid_iff}), and $\operatorname{pk}_b(p)$ is the number of peaks of $p$, plus $1$ if $b=1$ and $p$ begins with $D$ (\decl{peaksFrom_eq}); the flag $b$ records that the step before $p$ was a $U$. The count \decl{cnt} is
\[
  c(m,h,b,k)=\#\{p\in\Av_m:\ \operatorname{valid}_h(p),\ \operatorname{pk}_b(p)=k\}.
\]

\begin{lemma}[(a)--(c) in the statement file]\label{lem:cnt}
For all $m,h,k,n\in\N$ and $b\in\{0,1\}$:
\begin{enumerate}
\item[(a)] $T(n,k)=c(2n,0,0,k)$ \quad(\decl{T_eq_cnt});
\item[(b)] $c(m+2,h,b,k)=c(m+1,h+1,1,k)+\delta(m,h,b,k)+c(m,h,0,k)$, where $\delta=0$ if $h=0$; $\delta=c(m+1,h-1,0,k)$ if $h\ge1$, $b=0$; and $\delta=\ib{k\ge1}\,c(m+1,h-1,0,k-1)$ if $h\ge1$, $b=1$ \quad(\decl{cnt_step});
\item[(c)] $c(m,h,b,k)=0$ if $m<h$ or $m<k$ \quad(\decl{cnt_eq_zero});
\item[(d)] $c(0,0,b,k)=\ib{k=0}$, $c(1,1,0,k)=\ib{k=0}$ and $c(1,1,1,k)=\ib{k=1}$.
\end{enumerate}
\end{lemma}

\begin{proof}
(a) A word is a Schr\"oder path of semilength $n$ if and only if it has width $2n$ and $\operatorname{valid}_0$ holds, and $\operatorname{pk}_0$ is the number of peaks; as \texttt{words (2n)} has no repetitions, counting in the list and in the finset agree.
(b) $\Av_{m+2}=U\Av_{m+1}\sqcup D\Av_{m+1}\sqcup H\Av_m$, and the recursions above give the contribution of each part; a word $Dq$ is invalid from height $0$.
(c) A valid word from height $h$ has width $\ge h$ (\decl{valid_le}), and $\operatorname{pk}_b(p)\le|p|$ (\decl{peaksFrom_le}).
(d) $\Av_0=\{\varepsilon\}$ and $\Av_1=\{U,D\}$; from height $1$ only $D$ is valid, and $\operatorname{pk}_b(D)=b$ (proof file: \decl{cnt00}, \decl{W_base1f}, \decl{W_base1t}).
\end{proof}

\subsection{A bound on the number of peaks}

\begin{lemma}[\decl{peaks_bound}]\label{lem:bound}
If $\operatorname{valid}_h(p)$, then $2\operatorname{pk}_b(p)+h\le|p|+2b$.
\end{lemma}

\begin{proof}
Induction on $p$. If $p=\varepsilon$, validity forces $h=0$, and the left side is $0$. If $p=Uq$, then $\operatorname{valid}_{h+1}(q)$, and induction with flag $1$ gives $2\operatorname{pk}_1(q)+h+1\le|q|+2$, i.e.\ $2\operatorname{pk}_b(p)+h\le|q|+1=|p|$. If $p=Dq$, then $h\ge1$ and $\operatorname{valid}_{h-1}(q)$, so $2\operatorname{pk}_0(q)+h-1\le|q|$ and $2\operatorname{pk}_b(p)+h=2b+2\operatorname{pk}_0(q)+h\le |p|+2b$. If $p=Hq$, then $2\operatorname{pk}_b(p)+h=2\operatorname{pk}_0(q)+h\le|q|<|p|$.
\end{proof}

\begin{corollary}[\decl{cnt_peaks}]\label{cor:bound}
If $m<2k$ then $c(m,0,0,k)=0$. In particular $T(n,k)=0$ for $k>n$.
\end{corollary}

\subsection{Weighted sums}
For a weight $w\colon\N\to\Z$ put $W_w(m,h,b)=\sum_{k=0}^{m}w(k)\,c(m,h,b,k)$ (Lean \texttt{W w m h b}).
By Lemma~\ref{lem:cnt}(c) the range of summation may be replaced by $\{0,\dots,N-1\}$ for any $N\ge m+1$ (\decl{W_eq_range}); $W_w(m,h,b)=0$ if $m<h$ (\decl{W_zero_lt}); and $W_w$ is linear in $w$ (\decl{W_lin}). By Lemma~\ref{lem:cnt}(d), $W_w(0,0,b)=w(0)$, $W_w(1,1,0)=w(0)$ and $W_w(1,1,1)=w(1)$.

\begin{lemma}[\decl{W_step}]\label{lem:Wstep}
$W_w(m+2,h,b)=W_w(m+1,h+1,1)+W_w(m,h,0)+D_w(m,h,b)$, where $D_w(m,0,b)=0$, $D_w(m,h+1,0)=W_w(m+1,h,0)$, and $D_w(m,h+1,1)=W_{w(\cdot+1)}(m+1,h,0)$, with $w(\cdot+1)$ the weight $k\mapsto w(k+1)$.
\end{lemma}

\begin{proof}
Multiply Lemma~\ref{lem:cnt}(b) by $w(k)$ and sum over $0\le k\le m+2$, using \decl{W_eq_range}. For $h\ge1$, $b=1$, the sum coming from $\delta$ is $\sum_{k=1}^{m+2}w(k)c(m+1,h-1,0,k-1)=\sum_{j=0}^{m+1}w(j+1)c(m+1,h-1,0,j)$.
\end{proof}

\subsection{A closed-form principle}
Valid words from height $h$ have width $m\equiv h\pmod 2$; we write $m=h+2s$.

\begin{proposition}[\decl{closed_gen}]\label{prop:gen}
Let $w\colon\N\to\Z$, $F\colon\N\times\N\times\{0,1\}\to\Z$ and $L\colon\N\times\N\to\Z$ satisfy
\begin{enumerate}
\item[(i)] $F(0,0,b)=w(0)$ for both $b$, $F(1,0,0)=w(0)$ and $F(1,0,1)=w(1)$;
\item[(ii)] $W_{w(\cdot+1)}(m+1,h,0)=-W_w(m+1,h,0)+L(h,s)$ whenever $h+2s=m+1$;
\item[(iii)] for all $h,s,b$ with $h+2s\ge2$,
\begin{equation}\label{eq:rec}
\begin{aligned}F(h,s,b)={}&\ib{s\ge1}\bigl(F(h+1,s-1,1)+F(h,s-1,0)\bigr)\\
&+\ib{h\ge1}\cdot\begin{cases}F(h-1,s,0)&(b=0),\\ -F(h-1,s,0)+L(h-1,s)&(b=1).\end{cases}\end{aligned}
\end{equation}
\end{enumerate}
Then $W_w(m,h,b)=F(h,s,b)$ whenever $h+2s=m$.
\end{proposition}

\begin{proof}
Strong induction on $m$. If $m=0$ then $h=s=0$, and if $m=1$ then $h=1$, $s=0$; these cases are (i) with the base values of $W_w$. Let $m+2=h+2s$ and apply Lemma~\ref{lem:Wstep}. If $s=0$, then $m<h$, so $W_w(m+1,h+1,1)=W_w(m,h,0)=0$; if $s\ge1$, induction gives $W_w(m+1,h+1,1)=F(h+1,s-1,1)$ and $W_w(m,h,0)=F(h,s-1,0)$. If $h=0$ then $D_w(m,h,b)=0$. If $h\ge1$ then $(h-1)+2s=m+1$, so by induction $D_w(m,h,0)=F(h-1,s,0)$, and by (ii) and induction $D_w(m,h,1)=-F(h-1,s,0)+L(h-1,s)$. As $h+2s\ge2$, the total is $F(h,s,b)$ by \eqref{eq:rec}.
\end{proof}

\subsection{The three moments}
The weights $w_0(k)=(-1)^k$, $w_1(k)=(-1)^kk$, $w_2(k)=(-1)^kk(k-1)$ satisfy
\begin{equation}\label{eq:shift}
  w_0(k+1)=-w_0(k),\quad w_1(k+1)=-w_1(k)-w_0(k),\quad w_2(k+1)=-w_2(k)-2w_1(k);
\end{equation}
let
\begin{align*}
F_0(h,s,b)&=\ib{s=0}\,(1-2b\ib{h\ge1}),\qquad
F_1(h,s,b)=\begin{cases}-b\ib{h\ge1}&(s=0),\\ h-1&(s\ge1,\ b=1),\\ -(h+1)&(s\ge1,\ b=0),\end{cases}\\
F_2(h,s,b)&=\ib{s\ge1}\Bigl((h+1)(s-1)(s+h)+b\bigl(2h-2h(s-1)(s+h-1)\bigr)\Bigr).
\end{align*}

\begin{proposition}[\decl{W0_closed}, \decl{W1_closed}, \decl{W2_closed}]\label{prop:moments}
For $r=0,1,2$ and all $h,s,b$: $W_{w_r}(h+2s,h,b)=F_r(h,s,b)$.
\end{proposition}

\begin{proof}
Apply Proposition~\ref{prop:gen} with $(w,F,L)=(w_r,F_r,L_r)$, where $L_0=0$, $L_1(h,s)=-F_0(h,s,0)=-\ib{s=0}$ and $L_2(h,s)=-2F_1(h,s,0)=2(h+1)\ib{s\ge1}$, in the order $r=0,1,2$.

\emph{(ii).} By \eqref{eq:shift} and linearity, $W_{w_0(\cdot+1)}=-W_{w_0}$, $W_{w_1(\cdot+1)}=-W_{w_1}-W_{w_0}$ and $W_{w_2(\cdot+1)}=-W_{w_2}-2W_{w_1}$. For $r=1$ (resp.\ $r=2$), (ii) then follows from the already proved case $r=0$ (resp.\ $r=1$) at $(m+1,h,0)$.

\emph{(i).} $w_0(0)=1$, $w_0(1)=-1$, $w_1(0)=0$, $w_1(1)=-1$, $w_2(0)=w_2(1)=0$; these are $F_r(0,0,b)$, $F_r(1,0,0)$, $F_r(1,0,1)$.

\emph{(iii).} Let $h+2s\ge2$, let $R_r$ be the right side of \eqref{eq:rec} and $E_r=\ib{s\ge1}(F_r(h+1,s-1,1)+F_r(h,s-1,0))$ its first part.
If $s=0$, then $h\ge2$ and $E_r=0$, and $R_0=1$ $(b=0)$, $-1+0$ $(b=1)$; $R_1=0$ $(b=0)$, $-0+L_1(h-1,0)=-1$ $(b=1)$; $R_2=0$; these are $F_r(h,0,b)$.
Let $s\ge1$. For $r=0$: $F_0(h,s,b)=0$, $E_0=-1+1=0$ if $s=1$ and $E_0=0$ if $s\ge2$, and the second part involves only $F_0(h-1,s,0)=0$ and $L_0=0$.
For $r=1$: $E_1=-1$ both for $s=1$ ($-1+0$) and for $s\ge2$ ($h-(h+1)$); for $h\ge1$ the second part is $F_1(h-1,s,0)=-h$ if $b=0$ and $h+L_1(h-1,s)=h$ if $b=1$. So $R_1=-1-h$ resp.\ $-1+h$, also for $h=0$, which is $F_1(h,s,b)$.
For $r=2$, $s=1$: $F_2(h,1,b)=2bh$, $E_2=0$, and for $h\ge1$ the second part is $F_2(h-1,1,0)=0$ if $b=0$ and $L_2(h-1,1)=2h$ if $b=1$; so $R_2=2bh$.
For $r=2$, $s\ge2$: put $t=s-1\ge1$. Expanding,
\[
E_2=(h+2)(t-1)(t+h+1)+2(h+1)-2(h+1)(t-1)(t+h)+(h+1)(t-1)(t+h)=t(t+2h+1).
\]
With $F_2(h-1,s,0)=ht(t+h)$ and $L_2(h-1,s)=2h$ we get $R_2=E_2+ht(t+h)$ if $b=0$ and $R_2=E_2-ht(t+h)+2h$ if $b=1$ (for $h=0$ both equal $E_2$). Since $F_2(h,s,b)=(h+1)t(t+h+1)+b\,(2h-2ht(t+h))$, both cases reduce to $t(t+2h+1)+ht(t+h)=(h+1)t(t+h+1)$, and both sides equal $t(t+ht+h^2+2h+1)$.

In Lean, (iii) is proved by splitting into $s\in\{0,1,\ge2\}$, $h\in\{0,1,\ge2\}$ and $b$, and closing the $18$ goals for each $r$ by \texttt{simp}, followed by \texttt{omega} for $r=0$ and \texttt{ring\_nf} for $r=2$.
\end{proof}

\subsection{Proof of Theorem~\ref{thm:main}}
By Lemma~\ref{lem:cnt}(a) the sum in \eqref{eq:schulte} is $\sum_{k=0}^n(-1)^kc(2n,0,0,k)(n+1-k)^2$. By Corollary~\ref{cor:bound}, $c(2n,0,0,k)=0$ for $n<k\le2n$, so we may sum over $0\le k\le 2n$. Since $(n+1-k)^2=(n+1)^2-(2n+1)k+k(k-1)$, the sum equals
\[
  (n+1)^2\,W_{w_0}(2n,0,0)-(2n+1)\,W_{w_1}(2n,0,0)+W_{w_2}(2n,0,0).
\]
By Proposition~\ref{prop:moments} with $h=0$, $s=n$, $b=0$, the moments are $F_0(0,n,0)=\ib{n=0}$, $F_1(0,n,0)=-\ib{n\ge1}$ and $F_2(0,n,0)=\ib{n\ge1}(n-1)n$. For $n=0$ the sum is $1$; for $n\ge1$ it is $(2n+1)+n(n-1)=n^2+n+1$. \qed

\begin{remark}[the Narayana route; not formalized]\label{rem:narayana}
Let $N(n,j)$ be the number of Dyck paths of semilength $n$ with $j$ peaks, so $N(0,0)=1$, and for $n\ge1$, $N(n,0)=0$, $N(n,1)=1$ (only $U^nD^n$) and $N(n,2)=\binom n2$ \cite{A001263,A090181}. Replacing every $H$ of a Schr\"oder path by $UD$ gives a Dyck path whose peaks are the old peaks and the replaced steps; conversely, flattening any $j-k$ of the $j$ peaks of a Dyck path into $H$ steps creates no new peak. Hence $T(n,k)=\sum_jN(n,j)\binom jk$, i.e.\ $T_n(x)=\sum_jN(n,j)(1+x)^j$ (Harris's formula). Since $\sum_k(-1)^k\binom jk k^{\underline r}=(-1)^r r!\,\ib{j=r}$, we get $\sum_k(-1)^kk^{\underline r}T(n,k)=(-1)^rr!\,N(n,r)$, which for $n\ge1$ gives the moments $0$, $-1$, $n(n-1)$ found above.
\end{remark}

\section{Machine verification}\label{sec:lean}

\subsection*{Files and declarations}
The proof is in \path{lean/Math0/Frontier/OEIS/A060693.lean} in the math0 repository, \url{https://github.com/Foundation-0/math0}. It imports the frozen statement file \path{lean/Math0/Statements/OEIS/A060693.lean} and the audit commands in \path{lean/Math0/Audit.lean}. The main declaration is \decl{Math0.Frontier.A060693.result} \texttt{: statement}, followed by the command
\begin{lstlisting}
#math0_settles Math0.Frontier.A060693.result Math0.Statements.A060693.statement
\end{lstlisting}
All auxiliary declarations are private; their names are given with the corresponding results in Section~\ref{sec:proof}, and parts (a)--(c) of Lemma~\ref{lem:cnt} are proved in the statement file. Neither file uses \texttt{sorry}, \texttt{native\_decide}, custom axioms or \texttt{set\_option}. The proof file uses no \texttt{decide}; the only kernel evaluations are the data checks of the statement file, which the proof does not use.

\subsection*{Audit}
The command \texttt{\#math0\_settles} (in \path{lean/Math0/Audit.lean}, for statements that math0 formalizes itself) fails unless \decl{statement} is a \texttt{Prop}-valued definition without universe parameters; \decl{statement} and \decl{result} depend on no axioms besides \decl{propext}, \decl{Classical.choice} and \decl{Quot.sound} (this rejects \decl{sorryAx}, \decl{Lean.ofReduceBool} from \texttt{native\_decide}, and custom axioms); and the type of \decl{result} is syntactically \decl{statement} or $\neg$\,\decl{statement}, up to metadata, with no unfolding. Here it reports
\begin{lstlisting}[basicstyle=\ttfamily\scriptsize]
math0 audit: 'Math0.Statements.A060693.statement' depends on axioms [Classical.choice, Quot.sound, propext]
math0 audit: 'Math0.Frontier.A060693.result' depends on axioms [Classical.choice, Quot.sound, propext]
math0 statement check: 'Math0.Frontier.A060693.result' proves 'Math0.Statements.A060693.statement'
\end{lstlisting}

\subsection*{Toolchain and reproduction}
We use Lean~4 \cite{Lean4} \texttt{v4.33.1} and Mathlib \cite{Mathlib} \texttt{v4.33.1}, pinned in \path{lean/}; the repository also pins formal-conjectures \cite{FC} at commit \texttt{c77fdd4}, which this result does not use. To reproduce, run \path{scripts/verify.sh} from the repository root: it fetches the Mathlib cache, builds all math0 results with \texttt{lake build} (re-running every audit command) and checks that the pinned \texttt{formal\_conjectures} and \texttt{mathlib} checkouts are unmodified. The single file can be checked with \texttt{lake env lean Math0/Frontier/OEIS/A060693.lean} in \path{lean/}; on our machine this took under a minute, including loading Mathlib.

{\footnotesize
\begin{thebibliography}{99}

\bibitem{TRU}
ElonSG, \emph{prove: OEIS A060693 Schröder-path peak identity (Schulte 2017) --- v2: public peak law + Scribe repairs}, pull request \#7436 to \texttt{the-omega-institute/trureturing} (branch \texttt{dev}), merged September~13, 2026, \url{https://github.com/the-omega-institute/trureturing/pull/7436}; supersedes pull request \#7402; pre-registration issue \#7384.

\bibitem{FC}
M.~Firsching, P.~Lezeau, S.~Mercuri, M.~Z. Horv\'ath, Y.~Dillies, C.~S\"onne, E.~Wieser, F.~Zhang, T.~Hubert, B.~Ag\"uera y Arcas, and P.~Kohli, \emph{Formal Conjectures: An Open and Evolving Benchmark for Verified Discovery in Mathematics}, arXiv:2605.13171, 2026. Repository: \url{https://github.com/google-deepmind/formal-conjectures}.

\bibitem{Mathlib}
The mathlib Community, \emph{The Lean mathematical library}, in: Proceedings of the 9th ACM SIGPLAN International Conference on Certified Programs and Proofs (CPP 2020), ACM, 2020, 367--381.

\bibitem{Lean4}
L.~de Moura and S.~Ullrich, \emph{The Lean 4 theorem prover and programming language}, in: Automated Deduction -- CADE 28, Lecture Notes in Comput. Sci. \textbf{12699}, Springer, 2021, 625--635.

\bibitem{A001263}
The OEIS Foundation Inc., \emph{Triangle of Narayana numbers $T(n,k)=C(n-1,k-1)C(n,k-1)/k$ with $1\le k\le n$, read by rows}, entry A001263 (N.~J.~A. Sloane) in The On-Line Encyclopedia of Integer Sequences, \url{https://oeis.org/A001263}.

\bibitem{A060693}
The OEIS Foundation Inc., \emph{Triangle $(0\le k\le n)$ read by rows: $T(n,k)$ is the number of Schröder paths from $(0,0)$ to $(2n,0)$ having $k$ peaks}, entry A060693 (F.~Chapoton, 2001) in The On-Line Encyclopedia of Integer Sequences, \url{https://oeis.org/A060693}; revision \#147, October~4, 2026.

\bibitem{A090181}
The OEIS Foundation Inc., \emph{Triangle of Narayana (A001263) with $0\le k\le n$, read by rows}, entry A090181 (Ph.~Deléham, 2004) in The On-Line Encyclopedia of Integer Sequences, \url{https://oeis.org/A090181}.

\end{thebibliography}
}

\end{document}
